\documentclass[12pt,letterpaper,oneside,openany]{report}

% ------------------------------------------------------------
% 1. DOCUMENTO E IDIOMA
% ------------------------------------------------------------

\usepackage[utf8]{inputenc}
\usepackage[T1]{fontenc}
\usepackage[english,spanish,es-tabla]{babel}

\decimalpoint

% ------------------------------------------------------------
% 2. TIPOGRAFÍA
% ------------------------------------------------------------

\usepackage{newtxtext}
\usepackage{microtype}

% ------------------------------------------------------------
% 3. CONFIGURACIÓN DE PÁGINA - APA 7
% ------------------------------------------------------------

\usepackage[
    letterpaper,
    margin=1in,
    headheight=15pt
]{geometry}

\usepackage{setspace}
\doublespacing

\setlength{\parindent}{0.5in}
\setlength{\parskip}{0pt}
\setlength{\emergencystretch}{3em}

% ------------------------------------------------------------
% 4. MATEMÁTICAS Y FÍSICA
% ------------------------------------------------------------

\usepackage{amsmath}
\usepackage{amssymb}
\usepackage{mathtools}
\usepackage{bm}

% Unidades físicas
\usepackage{siunitx}

\sisetup{
    output-decimal-marker = {.},
    group-separator = {\,},
    group-minimum-digits = 5,
    per-mode = symbol,
    separate-uncertainty = true
}

% ------------------------------------------------------------
% 5. FIGURAS
% ------------------------------------------------------------

\usepackage{graphicx}
\usepackage{float}
\usepackage{caption}
\usepackage{xcolor}

% Directorios de imágenes
\graphicspath{
    {imagenes/}
    {imagenes/figuras/}
}

% Configuración de títulos de figuras y tablas

\captionsetup{
    font=normalsize,
    labelfont=bf,
    textfont=it,
    labelsep=newline,
    justification=raggedright,
    singlelinecheck=false,
    skip=0.25\baselineskip
}

% ------------------------------------------------------------
% 6. TABLAS
% ------------------------------------------------------------

\usepackage{booktabs}
\usepackage{array}

% ------------------------------------------------------------
% 7. CÓDIGO PYTHON
% ------------------------------------------------------------

\usepackage{listings}
\usepackage{upquote}

\lstdefinestyle{pythonAPA}{
    language=Python,
    basicstyle=\ttfamily\small,
    numbers=left,
    numberstyle=\tiny,
    numbersep=8pt,
    stepnumber=1,
    showstringspaces=false,
    breaklines=true,
    breakatwhitespace=false,
    keepspaces=true,
    columns=fullflexible,
    tabsize=4,
    frame=single,
    rulecolor=\color{black},
    captionpos=b,
    inputencoding=utf8,
    literate=
        {á}{{\'a}}1
        {é}{{\'e}}1
        {í}{{\'i}}1
        {ó}{{\'o}}1
        {ú}{{\'u}}1
        {Á}{{\'A}}1
        {É}{{\'E}}1
        {Í}{{\'I}}1
        {Ó}{{\'O}}1
        {Ú}{{\'U}}1
        {ñ}{{\~n}}1
        {Ñ}{{\~N}}1
        {ü}{{\"u}}1
        {Ü}{{\"U}}1
        {×}{{$\times$}}1
        {°}{{$^\circ$}}1
        {→}{{$\rightarrow$}}1
        {←}{{$\leftarrow$}}1
}

% ------------------------------------------------------------
% 8. LISTAS
% ------------------------------------------------------------

\usepackage{enumitem}

\setlist{
    topsep=0pt,
    itemsep=0pt,
    parsep=0pt,
    partopsep=0pt,
    leftmargin=0.5in
}

% ------------------------------------------------------------
% 9. ENCABEZADO Y NUMERACIÓN DE PÁGINAS
% ------------------------------------------------------------

\usepackage{fancyhdr}

\pagestyle{fancy}
\fancyhf{}

\fancyhead[R]{\thepage}

\renewcommand{\headrulewidth}{0pt}
\renewcommand{\footrulewidth}{0pt}

\fancypagestyle{plain}{
    \fancyhf{}
    \fancyhead[R]{\thepage}
    \renewcommand{\headrulewidth}{0pt}
}

% ------------------------------------------------------------
% 10. FORMATO DE TÍTULOS - APA 7
% ------------------------------------------------------------

\usepackage{titlesec}

% Nivel 1
\titleformat{\chapter}[block]
    {\normalfont\normalsize\bfseries\centering}
    {}{0pt}{}

\titlespacing*{\chapter}
    {0pt}{0pt}{\baselineskip}

% Nivel 2
\titleformat{\section}[block]
    {\normalfont\normalsize\bfseries}
    {}{0pt}{}

\titlespacing*{\section}
    {0pt}{\baselineskip}{0pt}

% Nivel 3
\titleformat{\subsection}[block]
    {\normalfont\normalsize\bfseries\itshape}
    {}{0pt}{}

\titlespacing*{\subsection}
    {0pt}{\baselineskip}{0pt}

% Nivel 4
\titleformat{\subsubsection}[runin]
    {\normalfont\normalsize\bfseries}
    {}{0pt}{}[.]

\titlespacing*{\subsubsection}
    {0pt}{\baselineskip}{0.5em}

% Nivel 5
\titleformat{\paragraph}[runin]
    {\normalfont\normalsize\bfseries\itshape}
    {}{0pt}{}[.]

\titlespacing*{\paragraph}
    {0pt}{\baselineskip}{0.5em}

% APA: encabezados sin numeración
\setcounter{secnumdepth}{-1}
\setcounter{tocdepth}{2}

% ------------------------------------------------------------
% 11. BIBLIOGRAFÍA - APA 7
% ------------------------------------------------------------

\usepackage{csquotes}

\usepackage[
    backend=biber,
    style=apa,
    sorting=nyt,
    sortcites=true,
    uniquelist=false
]{biblatex}

\DeclareLanguageMapping{spanish}{spanish-apa}

% Sangría francesa
\setlength{\bibhang}{0.5in}
\setlength{\bibitemsep}{0pt}

% ------------------------------------------------------------
% 12. HIPERVÍNCULOS Y REFERENCIAS CRUZADAS
% ------------------------------------------------------------

\usepackage[hidelinks]{hyperref}

\usepackage[
    nameinlink,
    noabbrev,
    capitalise
]{cleveref}

\crefname{figure}{figura}{figuras}
\Crefname{figure}{Figura}{Figuras}

\crefname{table}{tabla}{tablas}
\Crefname{table}{Tabla}{Tablas}

\crefname{equation}{ecuación}{ecuaciones}
\Crefname{equation}{Ecuación}{Ecuaciones}

% ------------------------------------------------------------
% 13. DATOS INSTITUCIONALES
% ------------------------------------------------------------

\newcommand{\universidad}{Universidad de Nariño}

\newcommand{\facultad}{
    Facultad de Ciencias Exactas y Naturales
}

\newcommand{\programa}{Programa de Física}

\newcommand{\ciudad}{San Juan de Pasto, Nariño}

% ------------------------------------------------------------
% 14. DATOS DEL INFORME
% ------------------------------------------------------------

\newcommand{\tituloInforme}{
    Modelación computacional de un acelerador lineal
}

\newcommand{\autorInforme}{
    Erick García\\
    Gabriela Herrera\\
    Juan Esteban Muñoz
}

\newcommand{\docenteInforme}{
    Mg. Jaime Betancourth
}

\newcommand{\asignaturaInforme}{
    Herramientas Computacionales para la Física II
}

\newcommand{\fechaInforme}{\today}




\addbibresource{bibliografia.bib}

\begin{document}

\selectlanguage{spanish}

% ============================================================
% PORTADA
% ============================================================

\pagenumbering{roman}

\begin{titlepage}

% Numeración APA: número de página arriba a la derecha
\thispagestyle{fancy}

\begin{center}

\vspace*{0.35in}

% ------------------------------------------------------------
% LOGO UNIVERSIDAD DE NARIÑO
% ------------------------------------------------------------

\includegraphics[
    width=0.78\textwidth,
    keepaspectratio
]{imagenes/udenar.png}

\vspace{0.65in}

% ------------------------------------------------------------
% TÍTULO DEL INFORME
% ------------------------------------------------------------

{\large\bfseries \tituloInforme\par}

\vspace{0.85in}

% ------------------------------------------------------------
% AUTORES
% ------------------------------------------------------------

{\normalsize
    \autorInforme\par
}

\vspace{0.55in}

% ------------------------------------------------------------
% INFORMACIÓN ACADÉMICA
% ------------------------------------------------------------

{\normalsize \asignaturaInforme\par}

\vspace{0.15cm}

{\normalsize
    Docente: \docenteInforme\par
}
% ------------------------------------------------------------
% INFORMACIÓN INSTITUCIONAL
% ------------------------------------------------------------
\vfill
{\normalsize\bfseries \universidad\par}

\vspace{0.15cm}

{\normalsize \facultad\par}

{\normalsize \programa\par}

\vspace{0.05cm}

% ------------------------------------------------------------
% LUGAR Y FECHA
% ------------------------------------------------------------

% \vfill

{\normalsize \ciudad\par}

\vspace{0.15cm}

{\normalsize \fechaInforme\par}

\end{center}

\end{titlepage}

\clearpage


\clearpage



\pagestyle{empty}

\pagestyle{empty}

\chapter*{Resumen}
\thispagestyle{empty}

% Entrada en el índice sin número de página
\addtocontents{toc}{
    \protect\contentsline{chapter}{Resumen}{}{}
}

Se desarrolló un modelo computacional de un acelerador
lineal de radiofrecuencia con tubos de deriva, de estructura
tipo Alvarez, que opera a una frecuencia de
$352.2\,\mathrm{MHz}$ y permite estudiar la aceleración de
protones desde $3$ hasta $50\,\mathrm{MeV}$. El diseño se
estableció a partir de la condición de sincronismo
$L_n=\beta_n\lambda$ y de la ganancia energética por cavidad,
$\Delta E=qV\cos\phi$, con $qV=1.0\,\mathrm{MeV}$.
Para una fase síncrona de $\phi_s=-30^\circ$, se obtuvieron
55 celdas y una longitud total de $10.59\,\mathrm{m}$,
correspondientes a un gradiente medio de aceleración de
$4.44\,\mathrm{MeV/m}$. Posteriormente, se simuló un paquete
de 3000 protones con una desviación estándar inicial de fase
de $15.0^\circ$ y una dispersión relativa de energía de
\SI{0.5}{\percent}. Los resultados muestran que, para
$\phi_s=-30^\circ$, el \SI{99.3}{\percent} del paquete
permanece agrupado longitudinalmente, mientras que la
desviación estándar de fase disminuye de $15.0^\circ$
a $5.3^\circ$, en concordancia con el amortiguamiento
adiabático. En contraste, para la operación en la cresta
de radiofrecuencia ($\phi_s=0^\circ$), únicamente el
\SI{24.1}{\percent} de las partículas alcanza la energía
final de diseño. Asimismo, la comparación entre la
simulación y la aproximación lineal para las oscilaciones
de sincrotrón presenta un avance de fase de $25.7^\circ$
por celda, frente a los $26.4^\circ$ predichos
teóricamente. En cuanto a la sincronización, se encontró
que un error del \SI{2}{\percent} en la frecuencia de
radiofrecuencia permite conservar el \SI{94.5}{\percent}
de las partículas capturadas, mientras que un error del
\SI{4}{\percent} reduce esta fracción al
\SI{55.5}{\percent}. Finalmente, se observa que los
electrones con velocidades cercanas a la de la luz
($\beta\approx 1$) no presentan oscilaciones de fase
apreciables.

\vspace{6pt}

\noindent
\hspace*{0.5in}
\textit{Palabras clave:}
LINAC, acelerador lineal, estructura de Alvarez, estabilidad de fase,
dinámica longitudinal de haces, oscilaciones de sincrotrón,
simulación computacional.


% ------------------------------------------------------------
% ABSTRACT
% ------------------------------------------------------------

\clearpage

\begin{otherlanguage}{english}

\chapter*{Abstract}
\thispagestyle{empty}

% Entry in the table of contents without page number
\addtocontents{toc}{
    \protect\contentsline{chapter}{Abstract}{}{}
}

A computational model of a radio-frequency drift-tube
linear accelerator (DTL) with an Alvarez-type structure,
operating at a frequency of $352.2\,\mathrm{MHz}$, was
developed to study proton acceleration from $3$ to
$50\,\mathrm{MeV}$. The design was based on the synchronism
condition $L_n=\beta_n\lambda$ and the energy gain per
accelerating gap, $\Delta E=qV\cos\phi$, with
$qV=1.0\,\mathrm{MeV}$. For a synchronous phase of
$\phi_s=-30^\circ$, the resulting design required 55 cells
and a total length of $10.59\,\mathrm{m}$, corresponding to
an average accelerating gradient of $4.44\,\mathrm{MeV/m}$.
Subsequently, a bunch of 3,000 protons was simulated with
an initial phase standard deviation of $15.0^\circ$ and a
relative energy spread of \SI{0.5}{\percent}. The results
show that, for $\phi_s=-30^\circ$, \SI{99.3}{\percent}
of the particle bunch remains longitudinally captured,
while the phase standard deviation decreases from
$15.0^\circ$ to $5.3^\circ$, consistent with adiabatic
damping. In contrast, when operating at the
radio-frequency crest ($\phi_s=0^\circ$), only
\SI{24.1}{\percent} of the particles reach the final
design energy. Furthermore, comparison of the simulated
synchrotron phase oscillations with the linear approximation
yields a phase advance of $25.7^\circ$ per cell, compared
with the theoretical prediction of $26.4^\circ$. Regarding
synchronization, a \SI{2}{\percent} radio-frequency error
allows \SI{94.5}{\percent} of the particles to remain
captured, whereas a \SI{4}{\percent} error reduces this
fraction to \SI{55.5}{\percent}. Finally, electrons
traveling at velocities close to the speed of light
($\beta\approx 1$) exhibit no appreciable phase oscillations.

\vspace{6pt}

\noindent
\hspace*{0.5in}
\textit{Keywords:}
linear accelerator, Alvarez structure, phase stability,
longitudinal beam dynamics, synchrotron oscillations,
computational simulation.

\end{otherlanguage}


\clearpage


\pagestyle{fancy}

\tableofcontents

\clearpage

% ============================================================
% CUERPO DEL INFORME
% NUMERACIÓN ARÁBIGA DESDE 1
% ============================================================

\pagenumbering{arabic}

% CAPÍTULO 1
 \chapter{Introducción}

Un acelerador lineal de partículas (linac) es un dispositivo
que permite incrementar la energía cinética de partículas
cargadas mediante campos eléctricos de radiofrecuencia (RF),
dispuestos a lo largo de una estructura de aceleración.
En los aceleradores lineales de tubos de deriva, particularmente
en la estructura tipo Alvarez, las partículas atraviesan
sucesivamente regiones de aceleración denominadas huecos,
en las cuales el campo eléctrico longitudinal les transfiere
energía. Entre estos huecos se encuentran los tubos de deriva,
que protegen a las partículas del campo eléctrico oscilante
durante su tránsito. De esta manera, la variación temporal
de la fase del campo permite que las partículas alcancen
los siguientes huecos en una fase adecuada para continuar
su aceleración \parencite{wangler,wille}.

El funcionamiento de este tipo de aceleradores requiere
mantener el sincronismo entre el movimiento de las partículas
y la oscilación del campo de radiofrecuencia. Para ello,
la longitud de las celdas debe ajustarse a la velocidad
de la partícula síncrona. A medida que esta incrementa
su energía cinética, también aumenta su velocidad, por
lo que las celdas deben adquirir longitudes progresivamente
mayores para conservar la condición de sincronismo.
Este principio constituye una característica fundamental
de los aceleradores lineales de tubos de deriva.

Los aceleradores lineales tienen diversas aplicaciones
en física y tecnología. Se emplean como etapas de inyección
en instalaciones de aceleración de partículas, en equipos
de radioterapia que utilizan haces de electrones con energías
típicamente comprendidas entre
$6\,\mathrm{MeV}$ y $18\,\mathrm{MeV}$,
y en instalaciones que producen neutrones mediante procesos
de espalación. Estas aplicaciones evidencian la importancia
de comprender los mecanismos de transferencia de energía
y las condiciones dinámicas que permiten mantener
la estabilidad del haz durante el proceso de aceleración.

En este contexto, el presente trabajo desarrolla un modelo
computacional de un acelerador lineal de protones con
estructura tipo Alvarez, operando a una frecuencia de
$352.2\,\mathrm{MHz}$, para estudiar la aceleración desde
$3\,\mathrm{MeV}$ hasta $50\,\mathrm{MeV}$.
El estudio comprende el diseño de las longitudes de las
celdas para la partícula síncrona, el análisis de la dinámica
longitudinal de un paquete de partículas y la evaluación
de la estabilidad de fase. Asimismo, se examinan las
oscilaciones de sincrotrón, el efecto de las desviaciones
en la frecuencia de radiofrecuencia y el comportamiento
de los electrones en el régimen ultrarrelativista.
Mediante la implementación numérica en Python, se busca
analizar la evolución del sistema y contrastar los resultados
de la simulación con las predicciones del modelo teórico.


% % CAPÍTULO 2
 \chapter{Objetivos}

\section{Objetivo general}

Desarrollar un modelo computacional de un acelerador
lineal de protones con estructura tipo Alvarez, mediante
la implementación de un modelo de aceleración por
radiofrecuencia, con el propósito de estudiar el diseño
de las celdas y analizar la dinámica longitudinal y la
estabilidad de fase de un paquete de partículas.

\section{Objetivos específicos}

\begin{enumerate}

    \item Implementar el modelo de aceleración por cavidades
    de radiofrecuencia, considerando la ganancia de energía
    cinética por hueco y la energía acumulada después de
    múltiples etapas de aceleración:

    \begin{equation}
        \Delta T_i = qV_i\cos\phi_i,
    \end{equation}

    \begin{equation}
        T_N = T_0 + \sum_{i=1}^{N} qV_i\cos\phi_i.
    \end{equation}

    \item Determinar las longitudes de las celdas del
    acelerador a partir de la condición de sincronismo
    correspondiente a la partícula síncrona.

    \item Analizar la influencia de la fase síncrona y
    de los errores de sincronización de la radiofrecuencia
    sobre la dinámica longitudinal y el comportamiento
    de un paquete de partículas.

    \item Contrastar los resultados de las simulaciones
    de las oscilaciones de sincrotrón con las predicciones
    de la teoría lineal y comparar el comportamiento
    longitudinal de protones y electrones.

\end{enumerate}



% % CAPÍTULO 3
% \input{capitulos/04_fundamento_teorico.tex}

% % CAPÍTULO 4
 \chapter{Método numérico}

\section{Parámetros y condiciones iniciales}

Para desarrollar las simulaciones numéricas se establecieron los parámetros físicos y computacionales del acelerador, definidos en el archivo \texttt{parametros.json}. El sistema considerado corresponde a un acelerador lineal de protones con estructura de tubos de deriva tipo Alvarez, diseñado para incrementar la energía cinética de las partículas desde un valor inicial de $3\,\mathrm{MeV}$ hasta una energía objetivo de $50\,\mathrm{MeV}$.

La frecuencia nominal del campo de radiofrecuencia se establece en
$f_{\mathrm{RF}}=352.2\,\mathrm{MHz}$, mientras que la amplitud de la tensión de aceleración por brecha corresponde a
$qV_{\mathrm{gap}}=1.0\,\mathrm{MeV}$.
La fase síncrona de referencia se fija en
$\phi_s=-30^\circ$.

La ganancia energética de la partícula síncrona en cada brecha está determinada por

\begin{equation}
    \Delta T_s=qV_{\mathrm{gap}}\cos(\phi_s).
\end{equation}

Para los parámetros seleccionados, se obtiene

\begin{equation}
    \Delta T_s
    =1.0\cos(-30^\circ)
    \approx 0.866\,\mathrm{MeV}.
\end{equation}

La longitud de onda asociada a la frecuencia de operación se calcula mediante

\begin{equation}
    \lambda=\frac{c}{f_{\mathrm{RF}}},
\end{equation}

donde $c$ representa la velocidad de la luz en el vacío. Para la frecuencia considerada se obtiene

\begin{equation}
    \lambda\approx0.8512\,\mathrm{m}.
\end{equation}

Para estudiar la dinámica longitudinal se considera un paquete de
$N_p=3000$ protones. Las condiciones iniciales se generan introduciendo dispersiones aleatorias en la fase y en la energía cinética respecto a la partícula síncrona.

La desviación inicial de fase se modela mediante una distribución normal de media cero y desviación estándar $\sigma_\phi=15^\circ$:

\begin{equation}
    \delta\phi_{i,0}\sim
    \mathcal{N}(0,\sigma_\phi^2).
\end{equation}

Por su parte, la desviación inicial de energía se define mediante

\begin{equation}
    \delta T_{i,0}\sim
    \mathcal{N}(0,\sigma_T^2),
\end{equation}

con

\begin{equation}
    \sigma_T=0.005T_0.
\end{equation}

En consecuencia, las condiciones iniciales de cada partícula $i$ están dadas por

\begin{equation}
    \phi_{i,0}=\phi_s+\delta\phi_{i,0},
\end{equation}

\begin{equation}
    T_{i,0}=T_0+\delta T_{i,0}.
\end{equation}

Las variables aleatorias utilizadas para generar ambas distribuciones se consideran independientes. Se establece una semilla aleatoria de $20261919$ para garantizar la reproducibilidad de las condiciones iniciales y de los resultados estadísticos de las simulaciones.

Además del caso de referencia, se estudian las fases síncronas
$-60^\circ$, $-15^\circ$, $0^\circ$ y $+30^\circ$, con el propósito de analizar la influencia de la fase de operación sobre la dinámica longitudinal del paquete.

Para estudiar la sensibilidad del acelerador frente a variaciones de la frecuencia de radiofrecuencia, se consideran los errores relativos

\begin{equation}
    \epsilon_f=
    \{0,\;0.002,\;0.005,\;0.01,\;0.02,\;0.04\},
\end{equation}

correspondientes a variaciones de $0\%$, $0.2\%$, $0.5\%$,
$1\%$, $2\%$ y $4\%$, respectivamente.

Finalmente, se incorpora un caso de comparación con electrones, para el cual se establece una energía cinética inicial de
$10\,\mathrm{MeV}$ y una estructura de 60 celdas. Este caso permite estudiar la evolución de fase de partículas que se encuentran desde su inyección en un régimen altamente relativista.

\subsection{Discretización de la dinámica longitudinal}

La evolución longitudinal de las partículas se modela mediante un procedimiento discreto que considera cada brecha de aceleración como una etapa del proceso. De esta manera, la dinámica se representa mediante ecuaciones en diferencias que permiten determinar la energía cinética, la velocidad relativista y la fase de cada partícula después de atravesar sucesivamente las celdas del acelerador.

El modelo considera dos variables dinámicas principales: la energía cinética $T_{i,n}$ y la fase de radiofrecuencia $\phi_{i,n}$, donde el subíndice $i$ identifica la partícula y $n$ representa la etapa de aceleración.

\subsubsection*{Actualización de la energía cinética}

La ganancia de energía de una partícula depende de la fase con la que llega a la brecha de aceleración. Por tanto, la energía cinética después de atravesar la brecha se determina mediante

\begin{equation}
    T_{i,n+1}
    =
    T_{i,n}
    +
    qV_{\mathrm{gap}}\cos(\phi_{i,n}).
\end{equation}

Esta expresión se aplica individualmente a cada partícula. En consecuencia, las diferencias iniciales de fase producen ganancias energéticas distintas, modificando progresivamente la distribución de energía del paquete.

\subsubsection*{Cálculo de la velocidad relativista}

Una vez actualizada la energía cinética, se determina la velocidad de cada partícula mediante las relaciones relativistas.

El factor de Lorentz se calcula como

\begin{equation}
    \gamma_{i,n+1}
    =
    1+\frac{T_{i,n+1}}{mc^2},
\end{equation}

mientras que el parámetro de velocidad está dado por

\begin{equation}
    \beta_{i,n+1}
    =
    \sqrt{1-\frac{1}{\gamma_{i,n+1}^{2}}}.
\end{equation}

Por consiguiente, la velocidad correspondiente a la nueva energía cinética es

\begin{equation}
    v_{i,n+1}=\beta_{i,n+1}c.
\end{equation}

Este procedimiento permite incorporar los efectos relativistas asociados al incremento de energía durante el proceso de aceleración.

\subsubsection*{Actualización de la fase}

La evolución de la fase depende del tiempo que tarda cada partícula en recorrer la longitud de una celda. Considerando que la velocidad calculada después de atravesar la brecha se utiliza para describir el tránsito por toda la celda, el tiempo de recorrido se aproxima mediante

\begin{equation}
    \Delta t_{i,n}
    =
    \frac{L_n}{\beta_{i,n+1}c}.
\end{equation}

La fase de radiofrecuencia se actualiza de acuerdo con

\begin{equation}
    \boxed{
    \phi_{i,n+1}
    =
    \phi_{i,n}
    +
    \omega\Delta t_{i,n}
    -
    2\pi
    }
\end{equation}

donde $\omega=2\pi f_{\mathrm{RF}}$ representa la frecuencia angular nominal.

El término $2\pi$ se resta para expresar el cambio de fase respecto al avance de un período completo de radiofrecuencia entre brechas consecutivas.

Para la partícula síncrona, la condición de sincronismo establece que la longitud de cada celda satisface

\begin{equation}
    L_n=\beta_{s,n+1}\lambda,
\end{equation}

donde $\beta_{s,n+1}$ corresponde al parámetro de velocidad de la partícula síncrona después de recibir la ganancia energética en la brecha.

En consecuencia, su tiempo de tránsito es

\begin{equation}
    \Delta t_s
    =
    \frac{\beta_{s,n+1}\lambda}
    {\beta_{s,n+1}c}
    =
    \frac{1}{f_{\mathrm{RF}}}.
\end{equation}

Por tanto, el avance de fase durante el tránsito corresponde a $2\pi$, de modo que la fase relativa de la partícula síncrona permanece constante de una brecha a la siguiente.

\subsection{Diseño de la estructura aceleradora}

El diseño longitudinal del acelerador se realiza mediante un procedimiento iterativo que determina la longitud de cada celda a partir de la energía cinética de la partícula síncrona.

El propósito es garantizar que el tiempo de tránsito entre dos brechas consecutivas corresponda a un período de la señal de radiofrecuencia, conservando así la condición de sincronismo.

Inicialmente se establecen la energía cinética de entrada $T_0$, la energía objetivo $T_f$, la amplitud de tensión $qV_{\mathrm{gap}}$, la fase síncrona $\phi_s$, la masa en reposo $m$ y la longitud de onda $\lambda$.

\subsubsection*{Evolución energética de la partícula síncrona}

Se considera que la partícula de referencia atraviesa cada brecha con una fase constante $\phi_s$. Por tanto, su energía cinética evoluciona de acuerdo con

\begin{equation}
    T_{s,n+1}
    =
    T_{s,n}
    +
    qV_{\mathrm{gap}}\cos(\phi_s).
\end{equation}

Para una fase síncrona constante, la ganancia energética por brecha también permanece constante.

\subsubsection*{Determinación de las longitudes de celda}

Una vez calculada la energía cinética después de cada brecha, se determina el factor de Lorentz de la partícula síncrona:

\begin{equation}
    \gamma_{s,n+1}
    =
    1+\frac{T_{s,n+1}}{mc^2}.
\end{equation}

A partir de este resultado se obtiene el parámetro de velocidad:

\begin{equation}
    \beta_{s,n+1}
    =
    \sqrt{1-\frac{1}{\gamma_{s,n+1}^{2}}}.
\end{equation}

La condición de sincronismo requiere que el tiempo de tránsito entre brechas consecutivas sea igual al período de radiofrecuencia:

\begin{equation}
    \Delta t_s=\frac{1}{f_{\mathrm{RF}}}.
\end{equation}

Como la distancia recorrida corresponde a $L_n=v_{s,n+1}\Delta t_s$, se obtiene

\begin{equation}
    L_n=\frac{\beta_{s,n+1}c}{f_{\mathrm{RF}}}.
\end{equation}

Finalmente, utilizando $\lambda=c/f_{\mathrm{RF}}$, la longitud de cada celda queda determinada por

\begin{equation}
    \boxed{
    L_n=\beta_{s,n+1}\lambda
    }.
\end{equation}

De esta manera, las longitudes de las celdas aumentan conforme se incrementa la velocidad de la partícula síncrona.

\subsubsection*{Procedimiento iterativo de diseño}

El algoritmo comienza con la energía cinética inicial y repite el cálculo de la ganancia energética, la velocidad relativista y la longitud de celda hasta alcanzar o superar la energía objetivo:

\begin{equation}
    T_{s,n+1}\geq T_f.
\end{equation}

En cada iteración se almacenan la energía cinética y la longitud calculada. Se establece además un número máximo de iteraciones como condición de seguridad computacional.

Debido a que el número de brechas es discreto, la energía final obtenida puede superar ligeramente el valor objetivo establecido. Esta diferencia se conserva en el análisis de resultados y se identifica como consecuencia de la discretización del proceso de aceleración.

\subsection{Seguimiento del paquete de partículas}

Una vez determinadas las longitudes de las celdas, se realiza el seguimiento individual de las partículas del paquete. A diferencia del procedimiento de diseño, que considera exclusivamente la partícula síncrona, este algoritmo incorpora las diferencias iniciales de energía y fase.

Para cada partícula $i$, se aplican sucesivamente las ecuaciones de actualización de energía, velocidad relativista y fase descritas anteriormente.

El procedimiento se repite a lo largo de todas las celdas, almacenando las trayectorias de energía y fase para su posterior análisis.

\subsubsection*{Criterios de captura longitudinal}

Para evaluar la evolución del paquete respecto a la partícula síncrona, se definen las desviaciones de energía y fase:

\begin{equation}
    \Delta T_{i,n}
    =
    T_{i,n}-T_{s,n},
\end{equation}

\begin{equation}
    \Delta\phi_{i,n}
    =
    \phi_{i,n}-\phi_s.
\end{equation}

Al finalizar la simulación, se considera que una partícula satisface el criterio de agrupamiento longitudinal cuando

\begin{equation}
    |\Delta\phi_{i,N_c}|<60^\circ,
\end{equation}

donde $N_c$ representa el número total de celdas.

El porcentaje de partículas que cumple esta condición se calcula mediante

\begin{equation}
    P_{\mathrm{cap}}
    =
    \frac{N_{\mathrm{cap}}}{N_p}\times100\%,
\end{equation}

donde $N_{\mathrm{cap}}$ corresponde al número de partículas que satisface el criterio angular.

Adicionalmente, se evalúa la proximidad energética respecto a la partícula síncrona mediante

\begin{equation}
    |\Delta T_{i,N_c}|<1\,\mathrm{MeV}.
\end{equation}

Este segundo criterio permite determinar la fracción de partículas cuya energía final se encuentra dentro de la tolerancia establecida.

Ambos indicadores se analizan por separado, puesto que describen propiedades diferentes de la distribución longitudinal del paquete.

\subsection{Simulación de las oscilaciones de sincrotrón}

Las oscilaciones de sincrotrón describen el movimiento longitudinal de una partícula alrededor de la trayectoria de referencia. Para estudiar este comportamiento se introduce una pequeña perturbación inicial de fase respecto a la partícula síncrona.

Se considera una partícula cuya fase inicial está dada por

\begin{equation}
    \phi_0=\phi_s+3^\circ.
\end{equation}

La ganancia energética de esta partícula está determinada por

\begin{equation}
    T_{n+1}
    =
    T_n+
    qV_{\mathrm{gap}}
    \cos(\phi_s+\Delta\phi_n),
\end{equation}

donde $\Delta\phi_n=\phi_n-\phi_s$ representa la desviación de fase.

Para pequeñas desviaciones, $|\Delta\phi_n|\ll1$, se utiliza la aproximación de Taylor:

\begin{equation}
    \cos(\phi_s+\Delta\phi_n)
    \approx
    \cos\phi_s-\sin\phi_s\Delta\phi_n.
\end{equation}

Esta expresión muestra que la diferencia de ganancia energética entre la partícula perturbada y la síncrona depende, en primer orden, de su desviación de fase.

La diferencia energética modifica la velocidad de la partícula y, por tanto, su tiempo de tránsito entre brechas. Como consecuencia, se producen nuevas desviaciones de fase que, bajo condiciones de estabilidad longitudinal, generan oscilaciones alrededor de la trayectoria de referencia.

En la aproximación lineal, el avance de fase de sincrotrón por celda se estima mediante

\begin{equation}
    \boxed{
    \mu=
    \sqrt{
    \frac{2\pi qV_{\mathrm{gap}}(-\sin\phi_s)}
    {mc^2\beta_s^2\gamma_s^3}
    }
    }.
    \label{eq:mu}
\end{equation}

Los parámetros relativistas corresponden a la partícula síncrona y se evalúan en la etapa de aceleración considerada.

Para obtener una estimación numérica del avance de fase, se identifica la posición de cruces por cero consecutivos y alternados de $\Delta\phi_n$. Si $a$ y $b$ representan las posiciones de dos cruces consecutivos, medidas en número de celdas, el avance de fase simulado se estima mediante

\begin{equation}
    \mu_{\mathrm{sim}}=\frac{\pi}{b-a}.
\end{equation}

Finalmente, se comparan los valores obtenidos mediante el seguimiento numérico con la predicción de la teoría lineal. Esta comparación permite evaluar la concordancia del modelo discreto con el comportamiento esperado para pequeñas oscilaciones.

\subsection{Sensibilidad a errores en la frecuencia de radiofrecuencia}

Para estudiar la sensibilidad del acelerador frente a variaciones de la frecuencia de operación, se introduce un error relativo respecto a la frecuencia nominal:

\begin{equation}
    \epsilon_f=
    \frac{f_{\mathrm{real}}-f_{\mathrm{RF}}}
    {f_{\mathrm{RF}}}.
\end{equation}

Despejando la frecuencia real se obtiene

\begin{equation}
    f_{\mathrm{real}}
    =
    f_{\mathrm{RF}}(1+\epsilon_f).
\end{equation}

En consecuencia, la frecuencia angular utilizada durante el seguimiento se modifica de acuerdo con

\begin{equation}
    \omega_{\mathrm{real}}
    =
    2\pi f_{\mathrm{RF}}(1+\epsilon_f).
\end{equation}

La ecuación de actualización de fase queda entonces expresada como

\begin{equation}
    \phi_{i,n+1}
    =
    \phi_{i,n}
    +
    \omega_{\mathrm{real}}\Delta t_{i,n}
    -
    2\pi.
\end{equation}

Para cada valor de $\epsilon_f$ se realiza el seguimiento de las partículas utilizando las longitudes de celda determinadas con la frecuencia nominal de diseño.

Se registra la evolución de las fases y energías para evaluar la acumulación de desviaciones respecto a la condición de sincronismo original. Finalmente, se comparan las fracciones de partículas agrupadas y sus energías finales para analizar la sensibilidad del acelerador frente a errores de frecuencia.

\subsection{Simulación del comportamiento longitudinal de electrones}

Como caso adicional, se estudia el comportamiento longitudinal de electrones altamente relativistas. Debido a su menor masa en reposo, estas partículas alcanzan velocidades próximas a la velocidad de la luz con energías cinéticas considerablemente menores que las requeridas por los protones.

Para los electrones se considera una energía cinética inicial de
$T_0=10\,\mathrm{MeV}$. Su energía en reposo es

\begin{equation}
    m_ec^2\approx0.511\,\mathrm{MeV}.
\end{equation}

Por tanto, el factor de Lorentz inicial es

\begin{equation}
    \gamma_0
    =
    1+\frac{10}{0.511}
    \approx20.57.
\end{equation}

El correspondiente parámetro de velocidad resulta

\begin{equation}
    \beta_0
    =
    \sqrt{1-\frac{1}{\gamma_0^2}}
    \approx0.9988.
\end{equation}

Este resultado indica que los electrones se encuentran desde su inyección en un régimen altamente relativista.

Para este caso se considera una estructura de 60 celdas de longitud constante, aproximada por la longitud de onda de radiofrecuencia:

\begin{equation}
    L=\lambda=\frac{c}{f_{\mathrm{RF}}}.
\end{equation}

La actualización de fase se determina mediante

\begin{equation}
    \phi_{i,n+1}
    =
    \phi_{i,n}
    +
    \omega\frac{L}{\beta_{i,n+1}c}
    -
    2\pi.
\end{equation}

Considerando la frecuencia nominal y la longitud constante $L=\lambda$, se obtiene

\begin{equation}
    \boxed{
    \phi_{i,n+1}
    =
    \phi_{i,n}
    +
    2\pi
    \left(
    \frac{1}{\beta_{i,n+1}}-1
    \right)
    }.
\end{equation}

Esta expresión permite analizar la deriva de fase asociada a la diferencia entre la velocidad de los electrones y la velocidad de la luz.

Para estudiar la influencia de la fase inicial se consideran cinco condiciones de inyección:

\begin{equation}
    \phi_0\in
    \{-40^\circ,-20^\circ,0^\circ,20^\circ,40^\circ\}.
\end{equation}

Para cada condición se realiza el seguimiento de la energía y la fase a través de las 60 celdas. Los resultados permiten estudiar la evolución longitudinal de partículas altamente relativistas y establecer una comparación con el comportamiento observado para los protones.

% % CAPÍTULO 5
% \chapter{Metodología}

% % CAPÍTULO 6
% \chapter{Desarrollo Computacional}

% % CAPÍTULO 7
 \chapter{Resultados}

\section{Diseño de la estructura aceleradora}

Para una fase síncrona de $\phi_s=-30^\circ$, la partícula de referencia experimenta una ganancia energética de aproximadamente
$0.866\,\mathrm{MeV}$ por brecha de aceleración. A partir de la condición de sincronismo y del procedimiento iterativo de diseño, se obtiene una estructura de 55 celdas, con una longitud total de $10.59\,\mathrm{m}$.

Las longitudes de las celdas aumentan progresivamente desde aproximadamente $7.7\,\mathrm{cm}$ hasta $26.9\,\mathrm{cm}$, en correspondencia con el incremento de la velocidad relativista de los protones. Este comportamiento concuerda con la condición de sincronismo $L_n=\beta_{s,n+1}\lambda$, según la cual la longitud de cada celda depende de la velocidad de la partícula síncrona después de atravesar la brecha anterior.

La energía cinética final obtenida mediante las 55 etapas de aceleración es aproximadamente $50.63\,\mathrm{MeV}$, valor ligeramente superior a la energía objetivo de $50\,\mathrm{MeV}$. Esta diferencia se debe a que el proceso de aceleración se realiza mediante un número entero de brechas.

El gradiente energético nominal, calculado a partir del incremento objetivo de $47\,\mathrm{MeV}$ y la longitud total de la estructura, es aproximadamente $4.44\,\mathrm{MeV/m}$.

La figura~\ref{fig:fig01_diseno} presenta la evolución de la energía cinética de la partícula síncrona y las longitudes correspondientes a las celdas del acelerador.

\subsection{Influencia de la fase síncrona en el diseño}

La fase síncrona determina la ganancia energética por brecha y, simultáneamente, las condiciones de estabilidad longitudinal. Para evaluar su influencia se realizaron diseños independientes utilizando diferentes valores de fase.

Los resultados muestran que, para $\phi_s=-60^\circ$, se requieren 94 celdas, mientras que para $\phi_s=0^\circ$ se necesitan 47 celdas. Esta diferencia se explica por la dependencia de la ganancia energética respecto a $\cos(\phi_s)$.

En particular, las fases más alejadas de la cresta proporcionan una menor ganancia energética por brecha. Por tanto, se requiere un mayor número de etapas para alcanzar la misma energía objetivo.

\section{Estabilidad longitudinal del paquete de partículas}

La evolución longitudinal del paquete se estudia mediante su distribución en el espacio de fases $(\Delta\phi,\Delta T)$. Para la configuración de referencia, correspondiente a $\phi_s=-30^\circ$, las partículas presentan oscilaciones alrededor de la trayectoria de la partícula síncrona.

Al finalizar la aceleración, la dispersión de fase disminuye de $15.0^\circ$ a $5.3^\circ$, mientras que la dispersión de energía alcanza aproximadamente $282\,\mathrm{keV}$. Este comportamiento es consistente con el amortiguamiento adiabático asociado al incremento de la energía de las partículas durante el proceso de aceleración.

De acuerdo con el criterio de captura establecido, el \SI{99.3}{\percent} del paquete permanece dentro de la ventana angular de $|\Delta\phi|<60^\circ$ al finalizar la estructura.

La tabla~\ref{tab:estabilidad} presenta los resultados obtenidos para las diferentes fases síncronas consideradas.

\begin{table}[H]
\centering
\caption{Comportamiento longitudinal del paquete para diferentes fases síncronas.}
\label{tab:estabilidad}
\small
\begin{tabular}{lccccc}
\toprule
$\phi_s$ ($^\circ$) &
Celdas &
Agrupadas (\%) &
Energía final (\%) &
$\sigma_\phi$ ($^\circ$) &
$\sigma_T$ (keV) \\
\midrule
-60 & 94 & 100.0 & 99.6 & 7.4 & 155 \\
-30 & 55 & 99.3 & 99.2 & 5.3 & 282 \\
-15 & 49 & 88.9 & 88.9 & 5.0 & 122 \\
0   & 47 & 26.8 & 24.1 & 14.4 & 457 \\
+30 & 55 & 41.4 & 45.9 & 6.8 & 547 \\
\bottomrule
\end{tabular}

\smallskip
\begin{minipage}{0.95\linewidth}
\footnotesize
Se consideraron $3000$ protones con una dispersión inicial de fase de $15^\circ$ y una dispersión relativa inicial de energía de \SI{0.5}{\percent}. Los diseños de las celdas se determinaron independientemente para cada fase síncrona. La columna de partículas agrupadas corresponde al criterio $|\Delta\phi|<60^\circ$, mientras que la energía final se evalúa mediante la tolerancia energética establecida en el método numérico.
\end{minipage}
\end{table}

Los resultados evidencian diferencias importantes en el comportamiento longitudinal dependiendo de la fase síncrona. Para $\phi_s=-60^\circ$ se obtiene la mayor fracción de partículas agrupadas, con un $100.0\,\%$, aunque se requiere una estructura considerablemente más larga.

Para $\phi_s=-15^\circ$, la dispersión final de fase es de $5.0^\circ$, pero la fracción de partículas agrupadas disminuye al $88.9\,\%$. Por otra parte, en la cresta de radiofrecuencia ($\phi_s=0^\circ$), la ausencia de una fuerza restauradora longitudinal de primer orden se refleja en una menor capacidad de mantener el paquete concentrado alrededor de la fase de referencia.

Para $\phi_s=+30^\circ$, la fase síncrona de diseño se encuentra fuera del intervalo de estabilidad lineal para protones. No obstante, la simulación muestra que una fracción del paquete evoluciona hacia regiones de fase distintas de la trayectoria de referencia, alcanzando una fracción de energía final del $45.9\,\%$. Este resultado debe interpretarse como una evolución no lineal del paquete y no como estabilidad de la fase síncrona inicialmente seleccionada.

La figura~\ref{fig:fig02_espacio_fases} presenta la distribución longitudinal del paquete para las configuraciones de referencia y para la fase positiva.

\section{Oscilaciones de sincrotrón}

Para estudiar las oscilaciones longitudinales se realizó el seguimiento de una partícula con una desviación inicial de fase de $3^\circ$ respecto a la partícula síncrona.

La evolución de la desviación de fase muestra un comportamiento oscilatorio alrededor de la trayectoria de referencia. A partir de los cruces por cero consecutivos se obtiene un avance de fase numérico de aproximadamente $25.7^\circ$ por celda en la región analizada, mientras que la predicción teórica correspondiente es de $26.4^\circ$ por celda.

La diferencia entre ambos valores es pequeña y resulta compatible con las limitaciones de resolución del procedimiento numérico utilizado para estimar el avance de fase a partir de posiciones discretas.

Adicionalmente, el avance de fase teórico disminuye desde aproximadamente $41.4^\circ$ hasta $9.8^\circ$ por celda a medida que aumenta la energía cinética de la partícula síncrona. Esta tendencia se encuentra asociada al incremento del factor relativista $\beta^2\gamma^3$ presente en el denominador de la expresión teórica.

Por otra parte, la amplitud de la desviación de fase disminuye desde $3.0^\circ$ hasta aproximadamente $1.39^\circ$. Esta reducción es consistente con el amortiguamiento adiabático de las oscilaciones longitudinales durante la aceleración.

La evolución de la amplitud presenta un comportamiento del orden de la dependencia $(\beta\gamma)^{-3/4}$. Entre las energías inicial y final consideradas, esta relación predice un factor aproximado de $0.35$, mientras que la simulación proporciona un factor de $0.46$. La diferencia puede estar relacionada con la evolución discreta de la oscilación y con la fase en la que se evalúa la amplitud en la última celda.

La figura~\ref{fig:fig03_sincrotron} muestra la evolución de la desviación de fase y la comparación entre el avance de fase teórico y el obtenido numéricamente.

\section{Sensibilidad a errores de frecuencia y comparación con electrones}

\subsection{Efecto de los errores en la frecuencia de RF}

Para analizar la sensibilidad del acelerador frente a variaciones de frecuencia, se realizaron simulaciones utilizando diferentes valores del error relativo $\epsilon_f$, manteniendo las longitudes de las celdas correspondientes al diseño nominal.

Los resultados muestran que las variaciones de frecuencia producen desplazamientos acumulativos de fase durante el tránsito por la estructura. Para un error relativo del \SI{1}{\percent}, el desplazamiento medio reportado es de $5.1^\circ$ y la energía cinética media final alcanza $51.69\,\mathrm{MeV}$, conservándose una fracción de partículas agrupadas del \SI{98.1}{\percent}.

Cuando el error aumenta al \SI{2}{\percent}, la fracción de partículas agrupadas disminuye al \SI{94.5}{\percent}. Para un error del \SI{4}{\percent}, esta fracción se reduce al \SI{55.5}{\percent}, lo que evidencia una mayor sensibilidad del paquete frente a desviaciones de la frecuencia nominal.

La tabla~\ref{tab:errorf} resume los resultados obtenidos.

\begin{table}[H]
\centering
\caption{Efecto de las variaciones de frecuencia de radiofrecuencia sobre el paquete de protones.}
\label{tab:errorf}
\small
\begin{tabular}{lccc}
\toprule
Error de frecuencia (\%) &
Partículas agrupadas (\%) &
$\langle T\rangle$ final (MeV) &
Desplazamiento de fase ($^\circ$) \\
\midrule
0.0 & 99.3 & 50.64 & 0.6 \\
0.2 & 99.0 & 50.86 & 1.4 \\
0.5 & 99.0 & 51.18 & 2.9 \\
1.0 & 98.1 & 51.69 & 5.1 \\
2.0 & 94.5 & 52.66 & 7.4 \\
4.0 & 55.5 & 55.50 & 10.7 \\
\bottomrule
\end{tabular}
\end{table}

Se observa que el incremento del error de frecuencia modifica tanto la distribución longitudinal como la energía media final del paquete. En particular, las diferencias entre las energías obtenidas y el valor nominal de diseño evidencian que el sincronismo depende de la frecuencia de operación del sistema.

\subsection{Comportamiento longitudinal de electrones}

Finalmente, se estudió el comportamiento de electrones con una energía cinética inicial de $10\,\mathrm{MeV}$ en una estructura de 60 celdas de longitud constante, aproximada por $\lambda$.

Para esta configuración, los electrones presentan inicialmente un factor de Lorentz de aproximadamente $20.57$ y una velocidad normalizada de $\beta_0\approx0.9988$. Por tanto, desde la inyección se encuentran en un régimen altamente relativista.

Los resultados de la simulación muestran que, para una inyección en la cresta de radiofrecuencia, la energía cinética final alcanza aproximadamente $69.92\,\mathrm{MeV}$. La fase experimenta una deriva total de aproximadamente $3.6^\circ$ a lo largo de la estructura, lo que evidencia que la velocidad próxima a la de la luz reduce significativamente la variación de fase durante el tránsito por las celdas.

Para las demás fases de inyección se observan diferencias en la energía final, debido a la dependencia de la ganancia energética respecto a la fase de llegada. En particular, una inyección con una fase de $40^\circ$ produce una energía final de aproximadamente $53.55\,\mathrm{MeV}$, de acuerdo con la menor ganancia energética asociada al factor $\cos(40^\circ)$.

La figura~\ref{fig:fig04_estabilidad_electrones} presenta los resultados correspondientes a la sensibilidad de la estabilidad longitudinal, las variaciones de frecuencia y la evolución de la energía de los electrones para diferentes fases de inyección.


% % CAPÍTULO 8
 \chapter{Validación y discusión}

\section{Validación del modelo numérico}

La consistencia del modelo computacional se evalúa mediante la comparación de sus resultados con las expresiones analíticas utilizadas durante el diseño del acelerador y con las predicciones de la teoría lineal de estabilidad longitudinal.

En primer lugar, se verifica la evolución energética de la partícula síncrona. Para una ganancia constante por brecha, la energía cinética después de $N_c$ etapas está determinada por

\begin{equation}
    T_{N_c}
    =
    T_0+
    N_cqV_{\mathrm{gap}}\cos(\phi_s).
\end{equation}

La comparación de esta expresión con los resultados del algoritmo de diseño permite comprobar la correcta implementación de la actualización energética.

En segundo lugar, se verifica la condición de sincronismo mediante la relación

\begin{equation}
    L_n=\beta_{s,n+1}\lambda.
\end{equation}

Esta condición garantiza que la partícula síncrona complete un período de radiofrecuencia durante el tránsito entre brechas consecutivas, conservando su fase de referencia.

Adicionalmente, se compara el avance de fase de las oscilaciones de sincrotrón obtenido mediante el seguimiento numérico con la predicción de la aproximación lineal, dada por la ecuación~\eqref{eq:mu}. Los resultados presentan una concordancia entre ambas estimaciones, dentro de la resolución del procedimiento discreto utilizado para determinar el avance de fase.

Desde el punto de vista cualitativo, las simulaciones también reproducen la dependencia de la estabilidad longitudinal respecto a la fase síncrona. Para protones, las fases situadas en el intervalo $-90^\circ<\phi_s<0^\circ$ presentan un comportamiento estable en la aproximación lineal, mientras que la operación en la cresta carece de una fuerza restauradora longitudinal de primer orden.

Finalmente, la utilización de una semilla aleatoria fija permite reproducir las mismas condiciones iniciales en ejecuciones sucesivas del programa. Sin embargo, la reproducibilidad de una realización particular no sustituye el análisis de la variabilidad estadística entre diferentes realizaciones del paquete.

\section{Limitaciones del modelo}

El modelo computacional permite estudiar la dinámica longitudinal del acelerador mediante un conjunto reducido de ecuaciones, pero incorpora simplificaciones que deben considerarse al interpretar los resultados.

En primer lugar, las brechas de aceleración se representan mediante una ganancia energética instantánea. Por tanto, no se considera la extensión espacial de las regiones aceleradoras ni la variación del campo eléctrico durante el tránsito dentro de ellas. En consecuencia, tampoco se incorpora un factor de tiempo de tránsito asociado a la distribución espacial del campo.

En segundo lugar, se considera una amplitud de tensión constante para todas las brechas. No se incluyen variaciones de amplitud entre cavidades ni fluctuaciones temporales de la señal de radiofrecuencia.

Asimismo, el modelo se limita al movimiento longitudinal de las partículas. No se consideran los grados de libertad transversales, la dinámica de focalización del haz ni los efectos de carga espacial. En un acelerador real, la estabilidad longitudinal puede estar acoplada a la dinámica transversal, por lo que se requieren elementos adicionales de focalización para controlar la divergencia del haz.

Por otra parte, el seguimiento se realiza mediante una discretización por celdas, utilizando la velocidad calculada después de cada brecha para aproximar el tiempo de tránsito por la celda correspondiente. Esta aproximación simplifica la integración de las ecuaciones de movimiento, pero no reproduce todos los detalles de la dinámica dentro de las cavidades.

Finalmente, la simulación del paquete considera una distribución estadística inicial de fase y energía, pero no incorpora una descripción completa de la distribución espacial del haz ni de sus interacciones colectivas.

Estas simplificaciones delimitan el alcance del modelo y deben tenerse presentes al extrapolar los resultados numéricos al comportamiento de un acelerador físico.

\section{Discusión}

Los resultados obtenidos permiten analizar la relación entre la condición de sincronismo, la estabilidad longitudinal y los parámetros de diseño de un acelerador lineal de protones con estructura tipo Alvarez.

La selección de la fase síncrona constituye uno de los aspectos fundamentales del diseño. Para $\phi_s=-30^\circ$, el acelerador alcanza una ganancia energética de aproximadamente $0.866\,\mathrm{MeV}$ por brecha y requiere 55 celdas para alcanzar la energía objetivo. Esta configuración proporciona una región de estabilidad longitudinal que permite mantener agrupada una fracción elevada del paquete de partículas.

Al modificar la fase síncrona se observa un compromiso entre la ganancia energética por brecha y las condiciones de estabilidad longitudinal. Las fases más alejadas de la cresta presentan una menor ganancia energética y, por tanto, requieren un mayor número de celdas. Sin embargo, también modifican la región longitudinal estable y la respuesta del paquete frente a las desviaciones iniciales de fase y energía.

La comparación entre las distintas fases síncronas muestra que la configuración de $-60^\circ$ proporciona la mayor fracción de partículas agrupadas entre los casos estudiados, aunque requiere una estructura considerablemente más larga que la configuración de $-30^\circ$. Por tanto, la elección de la fase síncrona no depende exclusivamente de maximizar la estabilidad, sino de establecer un compromiso entre la eficiencia energética y el comportamiento longitudinal del haz.

El análisis de las oscilaciones de sincrotrón permite identificar el mecanismo mediante el cual las desviaciones de fase y energía se acoplan durante el proceso de aceleración. La concordancia entre el avance de fase obtenido numéricamente y la predicción de la teoría lineal respalda la consistencia del modelo para pequeñas perturbaciones. Asimismo, la reducción de la amplitud de las oscilaciones es compatible con el amortiguamiento adiabático asociado al incremento de la energía relativista de las partículas.

Por otra parte, las simulaciones con errores de frecuencia muestran que el sincronismo es sensible a las variaciones de la señal de radiofrecuencia. A medida que aumenta el error relativo, se modifican las fases de llegada a las brechas y, en consecuencia, las ganancias energéticas experimentadas por las partículas. Para los parámetros considerados, los errores de hasta el $2\,\%$ conservan una fracción elevada del paquete dentro del criterio de agrupamiento establecido, mientras que un error del $4\,\%$ produce una reducción considerable de dicha fracción.

Finalmente, la comparación con electrones permite identificar diferencias importantes asociadas al régimen relativista. Debido a su menor masa en reposo, los electrones alcanzan velocidades próximas a la de la luz desde energías cinéticas relativamente bajas. En consecuencia, la deriva de fase por celda se reduce considerablemente cuando se utiliza una estructura de longitud constante aproximada por $\lambda$.

No obstante, es importante distinguir este comportamiento de la estabilidad longitudinal propiamente dicha. La reducción de la deriva de fase está relacionada con la proximidad de la velocidad de las partículas a $c$, mientras que la existencia de una fuerza restauradora longitudinal depende de la fase síncrona. En particular, para una operación en la cresta de radiofrecuencia, dicha fuerza desaparece en primer orden, independientemente de que la partícula acelerada sea un protón o un electrón.

\section{Conclusiones}

A partir del modelo computacional desarrollado y de los resultados obtenidos, se establecen las siguientes conclusiones:

\begin{enumerate}

    \item Se implementó un modelo discreto de aceleración por radiofrecuencia para una estructura de tubos de deriva tipo Alvarez. Mediante la condición de sincronismo y las relaciones relativistas se determinaron las longitudes de las celdas necesarias para acelerar protones desde $3\,\mathrm{MeV}$ hasta una energía objetivo de $50\,\mathrm{MeV}$. Para una fase síncrona de $-30^\circ$, se obtuvo una estructura de 55 celdas y una longitud total de $10.59\,\mathrm{m}$.

    \item El análisis del paquete de 3000 protones mostró que la fase síncrona influye significativamente en la estabilidad longitudinal. Para $\phi_s=-30^\circ$, el $99.3\,\%$ de las partículas satisface el criterio de agrupamiento establecido, mientras que para la operación en la cresta de radiofrecuencia esta fracción disminuye considerablemente. Estos resultados evidencian la importancia de seleccionar adecuadamente la fase síncrona para mantener la concentración longitudinal del haz.

    \item La simulación de las oscilaciones de sincrotrón presentó concordancia con las predicciones de la teoría lineal para pequeñas perturbaciones. Se observó una disminución del avance de fase por celda y de la amplitud de las oscilaciones conforme aumenta la energía relativista de la partícula, comportamiento consistente con el amortiguamiento adiabático.

    \item El análisis de errores en la frecuencia de radiofrecuencia mostró que las desviaciones respecto a la frecuencia nominal modifican la evolución de fase y la energía final del paquete. Para los parámetros estudiados, se conservó una elevada fracción de partículas agrupadas ante errores de frecuencia de hasta el $2\,\%$, mientras que un error del $4\,\%$ produjo una reducción considerable de la captura longitudinal.

    \item La comparación con electrones permitió identificar la influencia del régimen ultrarrelativista sobre la evolución de fase. La proximidad de su velocidad a la velocidad de la luz reduce la deriva de fase durante el tránsito por celdas de longitud constante. Sin embargo, la ausencia de una fuerza restauradora longitudinal de primer orden en la operación de cresta está determinada por la fase síncrona y no exclusivamente por la naturaleza de la partícula.

\end{enumerate}



% ============================================================
% REFERENCIAS
% ============================================================

\clearpage

\printbibliography[
    heading=bibintoc,
    title={Referencias}
]

% ============================================================
% APÉNDICES
% ============================================================

\appendix

% APÉNDICE A: CÓDIGO PYTHON

\clearpage

\chapter*{Apéndice A}

\addcontentsline{toc}{chapter}{Apéndice A. Código fuente}

\begin{center}
    \bfseries Código fuente
\end{center}

\lstinputlisting[
    style=pythonAPA,
    caption={Código fuente del proyecto},
    label={lst:codigo-python}
]{Apendices/codigo.py}

% APÉNDICE B: DATOS COMPLEMENTARIOS

\input{Apendices/datos.tex}

\end{document}
